% ================================================================
% ==== FILE: content/m_spaces/m2.tex
% ================================================================

% ==============================
%  Ontology of Continua — Core
%  Meta-Space: M2
%  FULL MODULE — FINAL
% ==============================

\subsubsection{\texorpdfstring{$M_2$}{M_2} Übersicht}
\label{sec:m2-overview}

$M_2$ ist der Metaraum, der die Zulässigkeitsstruktur für die bereitstellt
zweidimensionales Kontinuum $K_2$.
Der Übergang in $M_2$ entspricht der Entstehung von:

\begin{itemize}
    \item räumliche Dimensionalität über $A_1$ hinaus,
    \item lokale geometrische Struktur,
    \item eine dynamische Metrik,
    \item Kausalkegel und zulässige Ausbreitungsgeschwindigkeiten,
    \item-Quantisierungsschwellen und minimale Cluster,
    \item die vollständige aktionsbasierte Dynamik von Kontinua.
\end{itemize}

Es ist der erste Metaraum, in dem Geometrie, Zeit und physikalische Ausbreitung gemeinsam zulässig sind.

\subsubsection*{1. Struktur von \texorpdfstring{$\Omega(M_2)$}{\Omega(M_2)}}

Der zulässige Konfigurationsraum in $M_2$ ist:

\[
\Omega(M_2)
=
\Big\{
\phi : X_2 \to V \ \big|\
\phi \in C^{0}(X_2,V),\
\partial_i\phi,\ \partial_t\phi \in H^1,\
g_{\mu\nu} \in C^{0}
\Big\}.
\]

Anforderungen:

\begin{enumerate}
    \item \textbf{Zweidimensionale Topologie.}
    $X_2$ muss eine 2-Mannigfaltigkeit mit lokalen Koordinatendiagrammen sein.
    Typische Beispiele: $\mathbb{R}^2$, ein Zylinder oder eine kompakte Fläche.

    \item \textbf{Lokale Differenzierbarkeit.}
    Die ersten Ableitungen $\partial_i\phi$ und $\partial_t\phi$ müssen im Sinne von Sobolev existieren.

    \item \textbf{Zulässiges Metrikfeld.}
    $M_2$ erlaubt eine dynamische Metrik $g_{\mu\nu}$, aber nur mit Signatur:
    \[
    (-,+,+),
    \]
    entsteht als Hesse-Funktion des strukturellen Spannungsfunktionals $T_2$.

    \item \textbf{Ortsbeschränkung.}
    Zulässige Interaktionen müssen lokal in $X_2$ sein.

    \item \textbf{Endliche Ausbreitung.}
    Es müssen kausale Kegel vorhanden sein (Θ\(_{\text{conn}}\) erzwingt die Zulässigkeit von Lichtkegeln).
\end{enumerate}

Diese Bedingungen definieren den zulässigen Bereich für jedes $K_2$-Kontinuum.

\subsubsection*{2. Zulässige Achsen in \texorpdfstring{$M_2$}{M_2}}

$M_2$ lässt genau zwei unabhängige Raumachsen zu:

\[
A_1,\ \ A_2,
\]
mit $A_1$ geerbt und $A_2$ aus der Dimensionsschwellenbedingung:
\[
T_1 \ge \Theta_{\text{dim}}.
\]

In $M_2$ sind keine zusätzlichen Achsen zulässig.
Höhere Achsen erfordern Übergänge in $M_3$ und höher.

\subsubsection*{3. Zulässige Potenziale und Energien}

Der Metaraum $M_2$ erzwingt Energiefunktionale der Form:

\[
E[\phi]
=
\int_{X_2}
\left(
\frac{1}{2} g^{ij} \partial_i \phi \partial_j \phi
+ V(\phi)
\right) \sqrt{|g|}\, d^2x.
\]

Zulässige Potentiale $P(M_2)$ müssen Folgendes erfüllen:

\begin{itemize}
    \item Koerzitivkraft zur Gewährleistung der Stabilität gegen Kollaps,
    \item Begrenztheit von unten,
    \item Differenzierbarkeit zur Generierung von Flüssen,
    \item nur lokale Abhängigkeit (keine nichtlokalen Potentiale erlaubt).
\end{itemize}

\subsubsection*{4. Schwellen und strukturelle Spannungen}

$M_2$ definiert:

\[
\Theta_2 = \inf \{ T_2(\phi, g_{\mu\nu}) : (\phi, g) \in \Omega(M_2) \},
\]

mit der Strukturspannung:

\[
T_2 = \int_{X_2}
\left(
g^{ij} \partial_i \phi \partial_j \phi
+ W(\phi)
+ \alpha R
\right)
\sqrt{|g|}\, d^2x,
\]

wo:

\begin{itemize}
    \item $W(\phi)$ erzwingt Einschränkungen für zulässige Konfigurationen,
    \item $R$ ist die Skalarkrümmung,
    \item $\alpha$ steuert die geometrische Steifigkeit.
\end{itemize}

Ein Hauptmerkmal von $M_2$:
\\textbf{die Metrik $g_{\mu\nu}$ ergibt sich aus der zweiten Variante von $T_2$},
Geometrie zu einem abgeleiteten Konzept und nicht zu einem primitiven machen.

\subsubsection*{5. Zulässige Flüsse und Ausbreitung}

Zulässige Flüsse:

\[
J_2 = (\partial_i \phi,\ \partial_t \phi,\ \nabla_i T_2,\ \Box_g \phi),
\]

mit dem durch $g$ definierten d'Alembertian.

Ausbreitungsbeschränkungen:

\begin{itemize}
    \item-Signale müssen innerhalb des durch $\Theta_{\text{conn}}$ bestimmten Kausalkegels liegen,
    \item Ausbreitungsgeschwindigkeit darf den durch $g_{\mu\nu}$ festgelegten zulässigen Grenzwert nicht überschreiten,
    \item-Flüsse müssen die Regelmäßigkeit wahren ($H^1$-Kompatibilität).
\end{itemize}

Diese kodieren die Entstehung physikalischer Kausalität in $K_2$.

\subsubsection*{6. Quantisierungsschwellen und minimale Cluster}

$M_2$ ist der erste Metaraum, der Folgendes erzwingt:

\[
\Theta_{\text{quant}} > 0.
\]

Dies ergibt:

\begin{itemize}
    \item eine minimale stabile Clustergröße $q_{\min}$,
    \item diskrete Operatoraktionen:
    \[
    a^\dagger |k\rangle,\quad a |k\rangle,
    \]
    da Änderungen in der Konnektivität $k$ messen,
    \item quantisierte Anregungen als topologisch stabile Konfigurationen in $\Omega(M_2)$.
\end{itemize}

Somit ist Quantisierung eine strukturelle Notwendigkeit und kein externes Axiom.

\subsubsection*{7. Zusammenbruch und Grenze von \texorpdfstring{$M_2$}{M_2}}

Eine Konfiguration berührt $\partial\Omega(M_2)$, wenn:

\begin{itemize}
    \item die Metrik degeneriert ($\det g \to 0$),
    \item-Lokalität ist verletzt,
    \item $T_2 \ge \Theta_2$,
    \item Krümmung divergiert ($|R| \to \infty$),
    \item zulässige Achsen sind nicht mehr messbar.
\end{itemize}

Das Überschreiten dieser Grenze erzwingt den Zusammenbruch eines beliebigen $K_2$-Kontinuums.

\subsubsection*{8. Beziehung zu \texorpdfstring{$K_1$}{K_1} und \texorpdfstring{$K_3$}{K_3}}

$M_2$ ermöglicht den Übergang $K_1 \to K_2$ über $\Psi_{1\to2}$, definiert durch:

\[
A_1 \mapsto (A_1, A_2),\quad
\Omega(K_1) \to \Omega(K_2),\quad
T_1 \to T_2,\quad
\text{dim} = 2.
\]

Der Übergang zu $K_3$ wird nur zulässig, wenn:

\[
A_3 \in A(M_3),\qquad
T_2 \ge \Theta_{\text{dim}},
\qquad
\Omega(K_3) \subseteq \Omega(M_3).
\]

Somit ist $M_2$ der Metaraum von:

\begin{itemize}
    \item Geometrie,
    \item Kausalität,
    \item Quantisierung,
    \item Dimensionsstabilität,
    \item-Weitergabe.
\end{itemize}

Es ist die strukturelle Grundlage für alle höheren Kontinua.
